\chapter{Ejercicios y experimentos numéricos}

\section{Enunciado de los ejercicios}

\begin{enumerate}
	\item

	      Sean $\forall l=1,\dotsc,d:X_{l}\sim\mathcal{N}\left(0,1\right)$ son i.i.d.
	      Defina la variable aleatoria $X\coloneqq\sum\limits_{l=1}^{d}X^{2}_{l}$.

	      \begin{enumerate}
		      \item

		            Muestre que
		            \begin{math}
			            \forall t\in\left(0,\frac{1}{2}\right):
			            \mathbb{E}\left[\exp\left(tX\right)\right]=
			            \dfrac{1}{{\left(\sqrt{1-2t}\right)}^{d}}
		            \end{math}.

		      \item

		            Muestre que
		            \begin{math}
			            \forall a\in\left(0,\infty\right):
			            \mathbb{P}\left[X\geq a\right]\leq
			            \inf\limits_{0<t<\frac{1}{2}}
			            \dfrac{\exp\left(-ta\right)}{{\left(\sqrt{1-2t}\right)}^{d}}
		            \end{math}.
	      \end{enumerate}

	\item

	      Use la pregunta anterior para dar una demostración
	      alternativa del Teorema~\ref{thm:johnson-lindenstrauss},
	      donde no se requiere del teorema del anillo gaussiano.

	\item

	      En el siguiente ejercicio trataremos con datos generados aleatoriamente
	      con la proyección~\eqref{eq:projection},
	      calculamos el error actual de las distancias y comparamos esto con la cota
	      del Teorema~\ref{thm:johnson-lindenstrauss}.
	      Para $n=300$ muestras independientes que siguen
	      $X\sim\mathcal{N}\left(0,1,\mathbb{R}^{1000}\right)$
	      la siguiente gráfica emerge para el error multiplicativo actual
	      (línea sólida) comparada a la cota $\varepsilon=\sqrt{\frac{48\log n}{k}}$
	      (línea discontinua) provisto por Johnson-Lindenstrauss.
	      \begin{figure}[ht!]
		      \centering
		      \includegraphics[width=.55\paperwidth]{ejercicio3}
	      \end{figure}

	      Sean $d>k$ dados.
	      \begin{enumerate}
		      \item

		            Implemente la proyección~\eqref{eq:projection} y pruébelo para $d$ pequeño
		            y $k$.

		      \item

		            Haga $d=1000$, genere $300$ puntos aleatorios en
		            $\mathbb{R}^{d}$, de acuerdo a la distribución de su
		            elección, proyecte estos vía
		            $T_{U\left(\omega\right)}$ hacia
		            $\mathbb{R}^{k}$ y calcular para diferentes valores
		            de $k$ el mayor cambio de distancia que se produce
		            \begin{equation*}
			            \varepsilon\coloneqq%_{\text{JL}}
			            \max\left\{
			            1-
			            \min_{x\neq y}
			            \frac{
				            \left\|
				            T_{U\left(\omega\right)}\left(x\right)-
				            T_{U\left(\omega\right)}\left(y\right)
				            \right\|
			            }{
				            \left\|
				            x-y
				            \right\|
			            },
			            \max_{x\neq y}
			            \frac{
				            \left\|
				            T_{U\left(\omega\right)}\left(x\right)-
				            T_{U\left(\omega\right)}\left(y\right)
				            \right\|
			            }{
				            \left\|
				            x-y
				            \right\|
			            }-
			            1
			            \right\}.
		            \end{equation*}
	      \end{enumerate}
\end{enumerate}

\section{Solución de los ejercicios}

\begin{enumerate}
	\item

	      \begin{enumerate}
		      \item

		            \phantom{m}\vspace*{-1.6\baselineskip}\setlength\belowdisplayshortskip{0pt}
		            \begin{align*}
			            \shortintertext{Sea $t\in\left(0,\frac{1}{2}\right)$ fijo. La esperanza es}
			            \mathbb{E}
			            \left[
				            \exp
				            \left(tX\right)
			            \right] & =
			            \mathbb{E}
			            \left[
				            \exp
				            \left(
				            t\sum_{l=1}^{d}X^{2}_{l}
				            \right)
				            \right]=
			            \prod_{l=1}^{d}
			            \mathbb{E}
			            \left[\exp\left(tX^{2}_{l}\right)\right].  \\
			                    & \stackrel{\operatorname{def}}{=}
			            \prod_{l=1}^{d}
			            \int_{\mathbb{R}}
			            \frac{1}{\sqrt{2\pi}}
			            \exp
			            \left(t\tau^{2}\right)
			            \exp
			            \Big(-\frac{\tau^{2}}{2}\Big)
			            \dl\tau.                                   \\
			                    & =
			            \prod_{l=1}^{d}
			            \int_{\mathbb{R}}
			            \frac{1}{\sqrt{2\pi}}
			            \exp
			            \left[-\frac{\tau^{2}}{2}\left(1-2t\right)\right]
			            \dl\tau.
			            \shortintertext{Realice la $u$-sustitución: $u=\tau\sqrt{1-2t}$.}
			            \mathbb{E}
			            \left[
				            \exp
				            \left(tX\right)
			            \right] & =
			            \prod_{l=1}^{d}
			            \frac{1}{\sqrt{1-2t}}
			            \left[
				            \int_{-\infty}^{\infty}
				            \frac{1}{\sqrt{2\pi}}
				            \exp\left(-\frac{u^{2}}{2}\right)
				            \dl u
			            \right].                                   \\
			                    & =
			            \prod_{l=1}^{d}
			            \frac{1}{\sqrt{1-2t}}=
			            \frac{1}{{\left(\sqrt{1-2t}\right)}^{d}}.
		            \end{align*}

		      \item\label{list:1b}

		            Sea $a\in\left(0,\infty\right)$ fijo y $t\in\left(0,\frac{1}{2}\right)$.
		            \begin{align*}
			            \mathbb{P}\left[X\geq a\right] & =
			            \mathbb{P}\left[tX\geq ta\right]=
			            \mathbb{P}\left[\exp\left(tX\right)\geq\exp\left(ta\right)\right].
			            \shortintertext{Por la desigualdad de Markov: $\mathbb{P}\left[X\geq a\right]\leq\frac{1}{a}\mathbb{E}\left[X\right]$.}
			            \mathbb{P}\left[X\geq a\right] & \leq
			            \exp\left(-ta\right)
			            \mathbb{E}\left[\exp\left(tX\right)\right].
			            \shortintertext{Pero, en el ítem anterior se calculó la esperanza.}
			            \mathbb{P}\left[X\geq a\right] & \leq
			            \frac{\exp\left(-ta\right)}{{\left(\sqrt{1-2t}\right)}^{d}}.
			            \shortintertext{Finalmente, al tomar el ínfimo en $\left(0,\frac{1}{2}\right)$.}
			            \mathbb{P}\left[X\geq a\right] & \leq
			            \inf_{0<t<\frac{1}{2}}
			            \frac{\exp\left(-ta\right)}{{\left(\sqrt{1-2t}\right)}^{d}}.
		            \end{align*}
	      \end{enumerate}

	\item

	      Sea $\left\{\symbf{x}_{1},\dotsc,\symbf{x}_{n}\right\}\subset\mathbb{R}^{d}$.
	      Para cada par $i\neq j$, definimos $\symbf{x} \coloneqq \symbf{x}_i - \symbf{x}_j$.
	      La distancia proyectada al cuadrado es
	      \begin{equation*}
		      {\left\|
			      T_{U\left(\omega\right)}\left(\symbf{x}_{i}\right)-
			      T_{U\left(\omega\right)}\left(\symbf{x}_{j}\right)
			      \right\|}^{2} =
		      \frac{1}{k} {\left\| U \symbf{x} \right\|}^2 =
		      \frac{{\left\| \symbf{x} \right\|}^2}{k} \sum_{l=1}^k {\left\langle U_l, \frac{\symbf{x}}{\|\symbf{x}\|} \right\rangle}^2.
	      \end{equation*}
	      Puesto que $U_l \sim \mathcal{N}(\symbf{0}_d, I_d)$, se tiene que $Z_l \coloneqq \langle U_l, \frac{\symbf{x}}{\|\symbf{x}\|} \rangle \sim \mathcal{N}(0, 1)$ son i.i.d.
	      Definiendo $X \coloneqq \sum_{l=1}^k Z_l^2$, tenemos $X \sim \chi^2_k$, que es la variable del ejercicio anterior.
	      Usando los resultados de~\ref{list:1b} con
	      $a=k\left(1+\varepsilon\right)$ y
	      $a=k\left(1-\varepsilon\right)$, y las aproximaciones de
	      Taylor para el logaritmo $\log\left(1+x\right)\approx x-\frac{x^{2}}{2}$, se obtiene:
	      \begin{equation*}
		      \mathbb{P}
		      \left[
			      \frac{X}{k}\geq
			      1+\varepsilon
			      \right]\leq
		      \exp\left(-k\frac{\varepsilon^{2}}{4}\right)\text{ y }
		      \mathbb{P}
		      \left[\frac{X}{k}\leq 1-\varepsilon\right]\leq
		      \exp\left(-k\frac{\varepsilon^{2}}{4} \right).
	      \end{equation*}
	      Así, para un par fijo $\left(i,j\right)$, la probabilidad de
	      que la distorsión supere $\varepsilon$ es
	      \begin{equation*}
		      \mathbb{P}
		      \left[
			      \left|
			      \frac{
				      {
						      \left\|
						      T_{U\left(\omega\right)}
						      \left(\symbf{x}_{i}\right)-
						      T_{U\left(\omega\right)}
						      \left(\symbf{x}_j\right)
						      \right\|}^{2}}{
				      {\left\|\symbf{x}_{i}-\symbf{x}_{j}\right\|}^{2}
			      }-1
			      \right|\geq
			      \varepsilon
			      \right]\leq 2
		      \exp\left(-k\frac{\varepsilon^2}{4}\right).
	      \end{equation*}
	      Aplicando la cota de la unión sobre los $\binom{n}{2}$ pares de puntos:
	      \begin{equation*}
		      \mathbb{P}
		      \left[
			      \exists i,j:
			      \left|
			      \frac{
				      {\left\|T_{U}\left(\symbf{x}_{i}\right)-T_{U}\left(\symbf{x}_{j}\right)\right\|}^{2}
			      }{
				      {\left\|\symbf{x}_{i}-\symbf{x}_{j}\right\|}^{2}
			      }-1
			      \right|\geq
			      \varepsilon
			      \right]\leq
		      \sum_{i<j}
		      2\exp\left(-k\frac{\varepsilon^{2}}{4}\right)\leq
		      n^{2}
		      \exp\left(-k\frac{\varepsilon^2}{4}\right).
	      \end{equation*}
	      Sustituyendo el valor de
				$k\geq\frac{48}{\varepsilon^{2}}\log n$ se obtiene
	      \begin{equation*}
		      n^{2}
		      \exp\left(-k\frac{\varepsilon^2}{4}\right)\leq
		      n^{2}
		      \exp\left(-\frac{48\log n}{4}\right)=
		      n^{2}\exp\left(-12\log n\right)=n^{-10}\leq
		      \frac{1}{n}.
	      \end{equation*}
	      Por lo tanto, la probabilidad de que todos los pares
	      preserven sus distancias dentro de un factor
	      $\left(1\pm\varepsilon\right)$ es al menos $1-\frac{1}{n}$.

	\item

	      \begin{enumerate}
		      \item

		            La implementación de la proyección $T_{U\left(\omega\right)}$ se muestra en el apéndice~\ref{sec:projection}.

		      \item

		            Se ha realizado la siguiente tabla a partir del programa que se muestra en el apéndice~\ref{sec:projection}.
		            \begin{listing}[ht!]
			            \tiny
			            \centering
			            \inputminted[frame=single,framesep=10pt]{text}{../python/johnson.txt}
			            \caption{Comparación del error empírico y de la cota teórica de Johnson-Lindenstrauss para diferentes valores de $k$.}
			            \label{lst:results3a}
		            \end{listing}
	      \end{enumerate}
\end{enumerate}

\section{Distorsión de las distancias proyectadas $\left(k=4295\right)$}

Además, los resultados del Teorema~\ref{thm:johnson-lindenstrauss}
son totalmente independientes del conjunto original de puntos
$\left\{\symbf{x}_{1},\dotsc,\symbf{x}_{n}\right\}\subset\mathbb{R}^{d}$
o de la distribución que los generó (si la tuviesen).
De hecho, únicamente son dependientes del tamaño del conjunto sobre
el que se quieran aplicar.
Por ejemplo, considerando un conjunto de $10$ puntos en dimensión $d=10^{6}$,
entonces, por el Teorema~\ref{thm:johnson-lindenstrauss} se puede
proyectar los datos a un espacio de dimensión $k$ menor tal que las
distancias entre los proyectados y las distancias entre los puntos
originales difieren en menos de $\varepsilon=0.15$ con una probabilidad
mayor que $1-\frac{1}{10}=0.99$.
La dimensión del espacio de proyección debería ser mayor a
$\frac{48}{\left(0.15\right)^{2}}\log 10=4912.182$, es decir,
bastaría elegir $k=4913$, lo cual supondría una reducción del
espacio original a otro de dimensión unas $203$ veces más pequeña,
conservando las distancias con fidelidad.
La dimensión de este nuevo espacio de proyección, pese a seguir
siendo alta, se reduce considerablemente haciendo más manejable el
conjunto de datos original.

Pese a que la proyección al espacio de dimensión $k$ en el ejemplo anterior
traslada los datos a una dimensión considerablemente inferior, esta aún se puede mejorar.
Teniendo en cuenta el trabajo~\cite{Dasgupta2003}, se consigue mejorar la acotación inferior de $k$
demostrando que basta elegir
\begin{equation}\label{eq:dasguptaestimation}
	k\geq
	4
	\Big(
	\frac{\varepsilon^{2}}{2}-
	\frac{\varepsilon^{3}}{3}
	\Big)
	\log n,
\end{equation}
para que se cumpla el Teorema~\ref{thm:johnson-lindenstrauss}.
Volviendo al ejemplo anterior y, usando esta nueva acotación de $k$,
bastaría elegir $k$ mayor o igual que $909.66$ lo cual consigue reducir la
dimensión de forma aún más considerable.

Para comprobar experimentalmente los resultados del
Teorema~\ref{thm:johnson-lindenstrauss} se ha simulado en Python
una muestra de tamaño $n=150$ de la distribución
$X\sim\operatorname{Unif}\big({\left[0,1\right]}^{d}\big)$
con $d=100000$.
Se comprueba la distorsión promedio entre las distancias proyectadas
y las originales para una muestra de $500$ pares distintos de observaciones.

Para ello se ha escogido $\varepsilon=0.1$ y $k=4295$ según~\eqref{eq:dasguptaestimation}
y se ha generado una matriz aleatoria de tamaño $k\times d$ con
entradas gaussianas.

Por tanto, se representa en la Figura~\ref{fig:distorsion} el cociente entre los $500$ pares de distancias
proyectadas y las originales en el que se comprueba que dicho cociente es aproximadamente igual a $1$,
como lo establece el Teorema~\ref{thm:johnson-lindenstrauss}.

\begin{figure}[ht!]
	\centering
	\includegraphics[width=.75\paperwidth]{figura21_python}
	\caption{Distorsión de $500$ pares de puntos de una muestra de
		$150$ observaciones de la distribución uniforme en
		${\left[0,1\right]}^{d}$ con $d=100000$.}
	\label{fig:distorsion}
\end{figure}

\section[Comparación del Teorema para distintos epsilon y k]{Comparación del Teorema para distintos $\varepsilon$ y $k$}

Para generalizar lo anteriormente realizado, se ha calculado el promedio de las observaciones de esa misma muestra
que cumplen el Teorema~\ref{thm:johnson-lindenstrauss} escogiendo una muestra representativa de pares de observaciones.
Además, se ha variado el $\varepsilon$ y se ha considerado el $k$ dado por el Teorema~\ref{thm:johnson-lindenstrauss}
y el dado en la expresión~\eqref{eq:dasguptaestimation} con el fin de probar diferentes situaciones.

Estos resultados se han recogido en la Tabla~\ref{tab:verification} donde se observa cómo en todos
los casos la estimación promedio del porcentaje de pares de puntos cuya
distorsión está próxima a $1$ con una tolerancia de $\varepsilon$ es del $100\%$,
lo cual confirma los resultados teóricos del Teorema~\ref{thm:johnson-lindenstrauss}
de forma experimental.
Hemos considerado el $k$ en el Teorema~\ref{thm:johnson-lindenstrauss}
para apreciar la diferencia en la dimensión de proyección de los datos
entre esta y la que obtiene por la expresión~\eqref{eq:dasguptaestimation},
teniendo en cuenta un $\varepsilon$ mayor para no sobreexceder el límite de memoria
cuando se construye las matrices de proyección.

\begin{table}[ht!]
	\centering
	\caption{Comprobación del Teorema~\ref{thm:johnson-lindenstrauss} para distintos $\varepsilon$ y $k$.}
	\begin{tabular}{ccccc}
		\hline
		$\varepsilon$ & $k$ (con Teorema~\ref{thm:johnson-lindenstrauss}) & \%    & $k$ (con cota~\eqref{eq:dasguptaestimation}) & \%    \\
		\hline
		$0.20$        & $6013$                                            & $100$ & $1157$                                       & $100$ \\
		\hline
		$0.17$        & $8323$                                            & $100$ & $1565$                                       & $100$ \\
		\hline
		$0.15$        & $10690$                                           & $100$ & $1980$                                       & $100$ \\
		\hline
	\end{tabular}
	\label{tab:verification}
\end{table}
